%!TEX root = ../template.tex
%%%%%%%%%%%%%%%%%%%%%%%%%%%%%%%%%%%%%%%%%%%%%%%%%%%%%%%%%%%%%%%%%%%
%% chapter1.tex
%%%%%%%%%%%%%%%%%%%%%%%%%%%%%%%%%%%%%%%%%%%%%%%%%%%%%%%%%%%%%%%%%%%

\typeout{NT FILE chapter1.tex}%

\chapter{Introduction}
\label{cha:introduction}

%%------------------------------------------------------------
%% SECTION: Context \& Description
%%------------------------------------------------------------
\section{Context \& Description}
\label{sec:context-description}

% TODO: Situate the research domain, scope, and main phenomena or questions you address. Give the reader enough background to understand what follows without yet arguing why your work is needed.


%%------------------------------------------------------------
%% SECTION: Institutional Context
%%------------------------------------------------------------
\section{Institutional Context}
\label{sec:institutional-context}

% TODO: Describe the host institution, programme, funding, industrial partnership, or regulatory setting that frames the work, if relevant to how the problem is posed or validated.


%%------------------------------------------------------------
%% SECTION: Motivation
%%------------------------------------------------------------
\section{Motivation}
\label{sec:motivation}

% TODO: Explain the gap, limitation, or opportunity that drives this dissertation. Connect context to the concrete problem you tackle.


%%------------------------------------------------------------
%% SECTION: Objectives
%%------------------------------------------------------------
\section{Objectives}
\label{sec:objectives}

% TODO: State the main goal and, if useful, specific sub-objectives or research questions. Prefer measurable or verifiable formulations where possible.


%%------------------------------------------------------------
%% SECTION: Expected Contribution
%%------------------------------------------------------------
\section{Expected Contribution}
\label{sec:expected-contribution}

% TODO: Summarise what this work is expected to add (theoretical, methodological, empirical, or practical) and for whom it matters.


%%------------------------------------------------------------
%% SECTION: Structure
%%------------------------------------------------------------
\section{Structure}
\label{sec:structure}

% TODO: Briefly describe how the remaining chapters are organised.
% Example:
% Chapter~\ref{cha:background} reviews the relevant background and related work.
% Chapter~\ref{cha:design} presents the proposed approach.
% Chapter~\ref{cha:implementation} describes the implementation details.
% Chapter~\ref{cha:evaluation} reports on the evaluation and results.
% Chapter~\ref{cha:conclusion} concludes and outlines future work.
